\clearpage
\clearpage
\begin{center}
{\Large\bfseries 文献目録}
\end{center}
\label{bibliographie}

\begin{enumerate}
\setlength{\itemsep}{0.35em}
\item 三変数双二次形式の形式系の構成について。\\
Sitz. Ber. d. Physikal.-mediz. Sozietät in Erlangen 39 (1907), S. 176--179
\item 三変数双二次形式の形式系の構成について。\\
J. Reine Angew. Math. 134 (1908), S. 23--90 および表 1 点
\item \(n\) 変数形式の不変式論について。\\
J. Ber. d. DMV 19 (1910), S. 101--104
\item \(n\) 変数形式の不変式論について。\\
J. Reine Angew. Math. 139 (1911), S. 118--154
\item 有理関数体。\\
J. Ber. d. DMV 22 (1913), S. 316--319
\item 体と有理関数系。\\
Math. Ann. 76 (1915), S. 161--196
\item 有限群の不変式の有限性定理。\\
Math. Ann. 77 (1916), S. 89--92
\item 任意個の基本形式からなる系の不変式の整有理表示について。\\
Math. Ann. 77 (1916), S. 93--102.（第 16 番も参照）
\item 整超越数からなる最も一般的な領域。\\
Math. Ann. 77 (1916), S. 103--128.（第 16 番も参照）
\item 同型写像の関数方程式。\\
Math. Ann. 77 (1916), S. 536--545
\item 指定された群をもつ方程式。\\
Math. Ann. 78 (1918), S. 221--229.（第 16 番も参照）
\item 任意の微分式の不変式。\\
Nachr. v. d. Ges. d. Wiss. zu Göttingen 1918, S. 37--44
\item 不変変分問題。\\
Nachr. v. d. Ges. d. Wiss. zu Göttingen 1918, S. 235--257
\item 一変数代数関数の算術理論と、他の諸理論および代数体論との関係。\\
J. Ber. d. DMV 28 (1919), S. 182--203
\item 二変数形式の整係数不変式系の有限性。\\
Nachr. v. d. Ges. d. Wiss. zu Göttingen 1919, S. 138--156
\item 形式論における級数展開について。\\
Math. Ann. 81 (1920), S. 25--30
\item W. Schmeidler と共著：非可換領域における加群、特に微分式および差分式からなるもの。\\
Math. Z. 8 (1920), S. 1--35
\item 戦没した K. Hentzelt の消去理論に関する論文について。\\
J. Ber. d. DMV 30 (1921), S. 48
\item 環領域におけるイデアル論。\\
Math. Ann. 83 (1921), S. 24--66
\item 絶対既約性の代数的判定基準。\\
Math. Ann. 85 (1922), S. 26--33
\item 形式変分法と微分不変式。\\
Encyklopädie d. math. Wiss. III, 3 (1922), S. 68--71（R. Weitzenböck, Differentialinvarianten 所収）
\item K. Hentzelt \(\dagger\) の論文の編集：多項式イデアルと終結式の理論について。\\
Math. Ann. 88 (1923), S. 53--79
\item 代数不変式と微分不変式。\\
J. Ber. d. DMV 32 (1923), S. 177--184
\item 消去理論と一般イデアル論。\\
Math. Ann. 90 (1923), S. 229--261
\item 消去理論とイデアル論。\\
J. Ber. d. DMV 33 (1924), S. 116--120
\item 代数体におけるイデアル論の抽象的構成。\\
J. Ber. d. DMV 33 (1924), S. 102
\item イデアル論における Hilbert 数。\\
J. Ber. d. DMV 34 (1925), S. 101
\item 群指標とイデアル論。\\
J. Ber. d. DMV 34 (1925), S. 144
\item 標数 \(p\) の有限線形群の不変式の有限性定理。\\
Nachr. v. d. Ges. d. Wiss. zu Göttingen 1926, S. 28--35
\item 代数的な数体および関数体におけるイデアル論の抽象的構成。\\
Math. Ann. 96 (1927), S. 26--61
\item 代数的数体または関数体の整環に対する判別式定理。\\
J. Reine Angew. Math. 157 (1927), S. 82--104
\item R. Brauer と共著：既約表現の最小分解体について。\\
Sitz. Ber. d. Preuß. Akad. d. Wiss. 1927, S. 221--228
\item 算術的観点における超複素量と表現論。\\
Atti Congresso Bologna 2 (1928), S. 71--73
\item[34.] 超複素量と表現論。\\
Math. Z. 30 (1929), S. 641--692
\item[35.] 整関数からなる極大領域について。\\
Rec. Soc. Math. Moscou 36 (1929), S. 65--72
\item[36.] イデアル微分と共役差積。\\
J. Ber. d. DMV 39 (1930), S. 17
\item[37.] 高次分岐をもたない体の正規基底。\\
J. Reine Angew. Math. 167 (1932), S. 147--152
\item[38.] R. Brauer、H. Hasse と共著：代数の理論における主定理の証明。\\
J. Reine Angew. Math. 167 (1932), S. 399--404
\item[39.] 可換代数および数論との関係における超複素系。\\
Verhandl. Intern. Math.-Kongreß Zürich 1 (1932), S. 189--194
\item[40.] 非可換代数。\\
Math. Z. 37 (1933), S. 514--541
\item[41.] 相対ガロア数体に対する主種定理。\\
Math. Ann. 108 (1933), S. 411--419
\item[42.] 分解する接合積とその極大整環。\\
Actualités scientifiques et industrielles 148 (1934)（15 頁）
\item[43.] イデアル微分と共役差積。\\
J. Reine Angew. Math. 188 (1950), S. 1--21
\end{enumerate}

\noindent 代数関数に関する Noether の基本定理のための必要十分な重複度条件：H. Kapferer\\
（E. Noether との共同補遺を付す）。\\
Math. Ann. 97 (1927), S. 559--567

\clearpage
\begin{center}
{\Large\bfseries Emmy Noether の短報および書評一覧}
\end{center}
\label{liste-der-kurzmitteilungen-und-buchbesprechungen-von-emmy-noether}

\begin{center}
\emph{Deutsche Mathematiker-Vereinigung 年報掲載の短報}
\end{center}

\begingroup
\setlength{\parindent}{0pt}
\setlength{\parskip}{0.15em}
任意の微分式の不変式について、27, Teil II, (1918), S. 28\par
不変変分問題、27, Teil II, (1918), S. 47\par
二変数整係数不変式の有限性、28, Teil II, (1919), S. 29\par
整係数多項式と冪級数について、28, Teil II, (1919), S. 29--30\par
加群論の分解定理、29, Teil II, (1920), S. 54\par
単因子と一般イデアル論、30, Teil II, (1921), S. 32\par
Lipschitz の公式に関する指摘、30, Teil II, (1921), S. 48\par
一般イデアル論における Hilbert の定理の類似、32, Teil II, (1923), S. 20--21\par
公理論の問題に対する群論的注意、33, Teil II, (1925), S. 21\par
任意の体における Galois 理論、33, Teil II, (1925), S. 120\par
イデアル論における Hilbert 数、34, Teil II, (1926), S. 101\par
イデアルの微分商と分岐理論、38, Teil II, (1929), S. 81\par
\endgroup

\begin{center}
\emph{Deutsche Mathematiker-Vereinigung 年報掲載の書評}
\end{center}

\begingroup
\setlength{\parindent}{0pt}
\setlength{\parskip}{0.35em}
E. Study『ベクトル解析に基づく線形変換の不変式論入門』第 1 部（Die Wissenschaft, Bd. 71）、Braunschweig 1923, Vieweg。\\
36, Teil II, (1927), S. 71\par
Leopold Kronecker 著作集、第 4 巻、K. Hensel 編、VI + 509 頁、Leipzig und Berlin 1929, B. G. Teubner。\\
40, Teil II, (1931), S. 11--12\par
Leopold Kronecker 著作集、第 5 巻、K. Hensel 編、527 頁、Leipzig und Berlin 1930, B. G. Teubner。\\
41, Teil II, (1932), S. 17\par
Leopold Kronecker 著作集、第 III.2 巻、K. Hensel 編、215 頁、Leipzig und Berlin 1931, B. G. Teubner。\\
42, Teil II, (1933), S. 38--39\par
D. E. Rutherford『Modular Invariants』Cambridge Tracts No. 27、Cambridge 1932、VII + 84 頁。\\
43, Teil II, (1934), S. 94--95\par
\endgroup

\clearpage
\begin{center}
\emph{Emmy Noether の協力のもとに成立した書籍}
\end{center}
\label{buecher-entstanden-unter-mitwirkung-von-emmy-noether}

\begingroup
\setlength{\parindent}{0pt}
\setlength{\parskip}{0.5em}
R. Dedekind『数学著作集 I--III』R. Fricke、Emmy Noether、Øystein Ore 編、Vieweg, Braunschweig 1930--1932\\
（第 I 巻：1930、第 II 巻：1931、第 III 巻：1932）\par
G. Cantor『R. Dedekind との往復書簡』Emmy Noether、J. Cavaillès 編、60 頁、Actualités scient. et ind. 581, Paris 1937\par
\endgroup

\endgroup
